\subsection{Ordinary types and bounded carriers}
\label{sec:big-weighted-addition}

We first solve the weighted addition equation when one input has the ordinary type of $\varphi$, or a measure dominated by the Monge--Amp\`ere measure of a bounded potential. In both cases first moments control the rooftops of the approximating solutions, and these rooftops give the common lower bound required by \cref{lem:big-truncated-addition-limit}. The solution lies below both inputs, so its measure is finite.

\Needspace{14\baselineskip}
\begin{theorem}\label{thm:big-weighted-addition}
Let $u,v\in\mathcal E_\varphi$ and $\lambda>0$. Suppose either $v\sim\varphi$, or there are a Hermitian form $\omega$, a bounded $\psi\in\PSH(X,\omega)$ and a finite constant $A>0$ such that
\[
 \mu_v\leq A\,\mathrm{MA}_\omega(\psi).
\]
Then there is a unique $a\in\mathcal E_\varphi$ satisfying
\[
 \mu_a=\mathrm{e}^{\lambda(a-u)}\mu_u+\mathrm{e}^{\lambda(a-v)}\mu_v.
\]
It satisfies $a\leq\min(u,v)$, and hence $\mu_a\leq\mu_u+\mu_v$. If both $u,v\sim\varphi$, then also $a\sim\varphi$.
\end{theorem}
\begin{proof}
The proof has four steps: two inputs of minimal singularity type, two inputs of the ordinary type of $\varphi$, one input of that type, and finally an input with a bounded carrier measure.

It suffices to prove the result for $\lambda=1$. Indeed, replace $(\theta,\varphi,u,v)$ by $(\lambda\theta,\lambda\varphi,\lambda u,\lambda v)$: envelopes and relative classes scale with the potentials, volumes and Monge--Amp\`ere measures scale by $\lambda^n$, and a bounded-carrier domination remains one after changing its constant. Dividing the resulting solution by $\lambda$ gives the desired equation. Similarly, a common downward translation of $u$, $v$ and the solution allows us to assume $u,v\leq\varphi\leq0$; by the envelope identity, these inequalities hold after a sufficiently large common translation. We restore the translation at the end.

\emph{Step 1: inputs of minimal singularity type.}
First take $\varphi=V$ and $u,v\sim V$. Truncating only the weights, define the finite nonzero measures
\[
 \tau_j=\mathrm{e}^{-\max(u,-j)}\mu_u+\mathrm{e}^{-\max(v,-j)}\mu_v.
\]
The functions $\max(u,-j)$ need not be $\theta$-plurisubharmonic; they occur only as weights. \Cref{thm:big-minimal-general-measures} gives $a_j\in\mathcal E_V$ with $\mu_{a_j}=\mathrm{e}^{a_j}\tau_j$.

Choose $D\geq\tfrac12\|u-v\|_\infty+n\log2$, where the difference is taken off the common poles, and put $w=(u+v)/2-D$. Positivity of mixed products gives
\[
 \mu_w\geq2^{-n}(\mu_u+\mu_v)\geq \mathrm{e}^w\tau_j.
\]
Exponential order in $\mathcal E_V$ (\cref{cor:big-relative-orders}) gives $a_j\geq w$. Since $\tau_j$ increases with $j$, the same order shows that $a_j$ decreases. The envelope identity gives $a_1\leq V+\sup_Xa_1$, so $a_1\sim V$ by the lower bound $w$. Consequently
\[
 \mu_{a_j}
 \leq \mathrm{e}^{a_1-u}\mu_u+\mathrm{e}^{a_1-v}\mu_v
 \leq C(\mu_u+\mu_v).
\]
\Needspace{8\baselineskip}
The decreasing limit $a\geq w$ is not identically $-\infty$ and has the ordinary type of $V$. Fixed-dominator continuity (\cref{thm:big-relative-fixed-continuity}) gives weak convergence of the left-hand sides, and the right-hand sides converge in total variation by dominated convergence with respect to $\mu_u+\mu_v$. Therefore
\[
 \mu_a=\mathrm{e}^{a-u}\mu_u+\mathrm{e}^{a-v}\mu_v.
\]
Comparing this equation with each input measure, exponential order now gives $a\leq u,v$. This comparison uses the limiting equation; the truncated weights alone would not give it.

\emph{Step 2: two inputs of the ordinary type of $\varphi$.}
Suppose $u,v\sim\varphi$, and put $u_j=\max(u,V-j)$ and $v_j=\max(v,V-j)$. Step 1 gives solutions $a_j\sim V$ with
\[
 \mu_{a_j}=\mathrm{e}^{a_j-u_j}\mu_{u_j}+\mathrm{e}^{a_j-v_j}\mu_{v_j},
 \qquad a_j\leq\min(u_j,v_j).
\]
The differences $u_j-v_j$ satisfy the same uniform bound as $u-v$. With $D$ as above, $w_j=(u_j+v_j)/2-D$ satisfies the reverse inequality for this equation, and exponential order gives
\[
 (u+v)/2-D\leq w_j\leq a_j\leq\min(u_j,v_j).
\]
The fixed lower bound $w=(u+v)/2-D$ has the ordinary type of $\varphi$, and hence lies in $\mathcal E_\varphi$. Extract an $L^1$ limit $a\not\equiv-\infty$; it satisfies $w\leq a\leq\min(u,v)$, so $a\sim\varphi$. \Cref{lem:big-truncated-addition-limit}, with $\sigma=V$ and lower bound $h=w$, proves convergence in capacity and the equation. The lemma applies although the approximating solutions have the type of $V$ and the limit has the type of $\varphi$.

\emph{Step 3: one input of the ordinary type of $\varphi$.}
Now let $u\in\mathcal E_\varphi$ and $v\sim\varphi$, and set $u_j=\max(u,\varphi-j)$. Step 2 gives $a_j\sim\varphi$ with
\[
 \mu_{a_j}=\mathrm{e}^{a_j-u_j}\mu_{u_j}+\mathrm{e}^{a_j-v}\mu_v,
 \qquad a_j\leq\min(u_j,v).
\]
We construct a common lower bound in $\mathcal E_\varphi$. Put $\delta=n\log2>0$ and
\[
 s_j=Q_2(a_j,u_j)\in\mathcal E_\varphi,
 \qquad C_j=\{s_j+u_j=2a_j\}.
\]
\Needspace{6\baselineskip}
By stretching and contact, $s_j\leq a_j$ and
\begin{equation}\label{eq:big-addition-stretch-contact}
 2^{-n}\mu_{s_j}+2^{-n}\mathbf1_{C_j}\mu_{u_j}
 \leq\mathbf1_{C_j}\mu_{a_j}.
\end{equation}
Put $D_j=\{s_j\leq u-2\delta\}$. On $D_j\cap C_j$ we have $a_j=(s_j+u_j)/2\leq u_j-\delta$, because $u\leq u_j$. Since $\mathrm{e}^{-\delta}=2^{-n}$, inserting the equation for $\mu_{a_j}$ into \eqref{eq:big-addition-stretch-contact} shows that the $\mu_{u_j}$ terms cancel there. Also $v\geq\varphi-C\geq u_j-C$ for a fixed constant $C$. Hence
\begin{equation}\label{eq:big-addition-ordinary-tail}
 \mathbf1_{D_j}\mu_{s_j}
 \leq2^n \mathrm{e}^C\mathbf1_{D_j\cap C_j}
                \mathrm{e}^{(s_j-u_j)/2}\mu_v.
\end{equation}
All measures being subtracted here are finite.

\Needspace{6\baselineskip}
For $N\geq1$ form the finite rooftops
\[
 R_N=P_\theta(s_1,\ldots,s_N),\qquad q_N=P_\theta(u,R_N).
\]
They belong to $\mathcal E_\varphi$ by the finite-rooftop proposition. On the set $\{R_N\leq u-2\delta\}$, partition the contact set according to the first index $i$ with $R_N=s_i$. Since $u_i\geq u_N$, \eqref{eq:big-addition-ordinary-tail} gives
\[
 \mathbf1_{\{R_N\leq u-2\delta\}}\mu_{R_N}
 \leq2^n \mathrm{e}^C \mathrm{e}^{(R_N-u_N)/2}\mu_v.
\]
\Needspace{7\baselineskip}
The contact pieces are disjoint, so there is no factor $N$. On $\{q_N\leq u-2\delta\}$ only the contact part $q_N=R_N$ contributes. Since $u\leq u_N$ and $t \mathrm{e}^{-t/2}\leq2/\mathrm{e}$ for $t\geq0$, we obtain
\[
 \int_X(u-q_N)\,\mu_{q_N}
 \leq 2\delta M+2^{n+1}\mathrm{e}^{C-1}\mu_v(X).
\]
The first term bounds the contribution of the complementary set, including any contact with $u$. The first-moment criterion \cref{prop:big-relative-first-moment}, with fixed reference $u$, implies that the decreasing sequence $q_N$ has a limit $h\in\mathcal E_\varphi$, not identically $-\infty$. For each fixed $j$ and every $N\geq j$, we have $q_N\leq s_j\leq a_j$, so $h\leq a_j$ for every $j$.

The $a_j$ now have a subsequence converging in $L^1$ to a limit not identically $-\infty$, bounded above by $u$ and $v$. \Cref{lem:big-truncated-addition-limit}, with reference $\varphi$, fixed second input $v$ and common lower bound $h$, proves the required equation in $\mathcal E_\varphi$.

\emph{Step 4: a bounded carrier measure for the second input.}
Finally let $u,v\in\mathcal E_\varphi$ and $\mu_v\leq A\nu$, with $\nu=\mathrm{MA}_\omega(\psi)$ and $\psi$ bounded. Again put $u_j=\max(u,\varphi-j)$. By Step 3, applied with the inputs interchanged, there are $a_j\in\mathcal E_\varphi$ satisfying the same equations and bounds as in Step 3.

Let $r=P_\theta(u,v)\in\mathcal E_\varphi$. On $\{a_j<r\}$ one has $r\leq u_j,v$, so the inequality $t \mathrm{e}^{-t}\leq1/\mathrm{e}$ gives
\[
 \int_X(r-a_j)_+\,\mu_{a_j}
 \leq \mathrm{e}^{-1}\bigl(\mu_{u_j}(X)+\mu_v(X)\bigr)\leq2M/\mathrm{e}.
\]
\Needspace{7\baselineskip}
Since $a_j\leq\varphi$, the one-sided first-moment criterion gives a uniform lower bound for $\sup_Xa_j$. By \cref{lem:bounded-carrier-integrability},
\[
 \sup_{\substack{z\in\PSH(X,\theta)\\\sup_Xz=0}}
       \int_X(-z)\,\mathrm{d}\nu<\infty.
\]
\Needspace{6\baselineskip}
Applying this to $a_j-\sup_Xa_j$ and using the lower bound for the suprema, we obtain
\begin{equation}\label{eq:big-addition-carrier-moment}
 \sup_j\int_X(u_j-a_j)\,\mu_v
 \leq\sup_j\int_X(-a_j)\,\mu_v<\infty.
\end{equation}

Keep $s_j=Q_2(a_j,u_j)$, $C_j$ and $D_j$ from Step 3. The same cancellation in \eqref{eq:big-addition-stretch-contact}, now using only $a_j\leq v$, gives
\[
 \mathbf1_{D_j}\mu_{s_j}
 \leq2^n\mathbf1_{D_j\cap C_j}\mathrm{e}^{a_j-v}\mu_v
 \leq2^n\mu_v.
\]
\Needspace{6\baselineskip}
Take one additional rooftop,
\[
 q_j=P_\theta(s_j,u-2\delta)\in\mathcal E_\varphi.
\]
On its contact set with $u-2\delta$ its measure is at most $\mu_u$; on its other contact part, where $q_j=s_j\leq u-2\delta$, the preceding inequality applies. Thus
\begin{equation}\label{eq:big-addition-fixed-roof-measure}
 \mu_{q_j}\leq\mu_u+2^n\mu_v.
\end{equation}
Moreover the integrand $u-2\delta-q_j$ vanishes on the first contact part. On the other part, the measure is carried by $C_j$, where $u_j-s_j=2(u_j-a_j)$. Therefore
\[
 \int_X(u-2\delta-q_j)\,\mu_{q_j}
 \leq2^{n+1}\int_X(u_j-a_j)\,\mu_v,
\]
which is uniformly bounded by \eqref{eq:big-addition-carrier-moment}. The first-moment criterion with reference $u-2\delta$ prevents collapse of the $q_j$. Extract $L^1$ limits of the $a_j$ and of the $q_j$, not identically $-\infty$. The fixed-dominator theorem applies to \eqref{eq:big-addition-fixed-roof-measure} and, after a further extraction, gives
\[
 h=P_\theta\left(\inf_jq_j\right)\in\mathcal E_\varphi.
\]
Since $q_j\leq s_j\leq a_j$, this is a common lower bound for the extracted solutions, and \cref{lem:big-truncated-addition-limit} gives the equation. Its truncation levels are the original indices of the extracted subsequence.

In every step the solution lies below $u$ and $v$, so its measure is bounded by $\mu_u+\mu_v$. Two solutions have the same exponentially weighted measure, so \cref{cor:big-relative-orders}, applied in both directions, gives uniqueness, even when that weighted measure has infinite mass. Restoring the common translation and the exponent finishes the proof.
\end{proof}
